% =============================================================================
% Annexe B : Configurations
% =============================================================================
\chapter{Configurations}
\label{annexe:configurations}

\section{Variables d'environnement (.env)}

\begin{lstlisting}[language=bash, caption={.env.example --- Variables de configuration des tests}, label={lst:env}]
# URLs de base
ADMIN_BASE_URL=http://localhost:4200
PARTNER_BASE_URL=http://localhost:4201
API_GATEWAY_URL=http://localhost:8080

# Comptes de test (crees par le script de seed)
ADMIN_EMAIL=admin@speedline-test.com
ADMIN_PASSWORD=TestAdmin123!
PARTNER_EMAIL=partner@speedline-test.com
PARTNER_PASSWORD=TestPartner123!
CUSTOMER_EMAIL=customer@speedline-test.com
CUSTOMER_PASSWORD=TestCustomer123!
COURIER_EMAIL=courier@speedline-test.com
COURIER_PASSWORD=TestCourier123!

# Base de donnees
DB_HOST=localhost
DB_PORT=5432
DB_USER=postgres
DB_PASSWORD=postgres123
AUTH_DB_NAME=speedline_auth
\end{lstlisting}

\section{Package.json --- Dépendances du robot}

\begin{lstlisting}[language=Java, caption={package.json --- Scripts et dépendances npm}, label={lst:package-json}]
{
  "name": "speedline-test-robot",
  "version": "1.0.0",
  "scripts": {
    "test": "npx playwright test",
    "test:admin": "npx playwright test --project=admin-panel",
    "test:partner": "npx playwright test --project=partner-dashboard",
    "test:api": "npx playwright test --project=api-tests",
    "test:smoke": "npx playwright test api/smoke.api.spec.ts",
    "test:security": "npx playwright test api/security.api.spec.ts",
    "test:visual": "npx playwright test visual-regression",
    "test:a11y": "npx playwright test accessibility",
    "report:allure": "npx allure generate allure-results --clean",
    "seed": "node ../seed/seed-test-data.js"
  },
  "devDependencies": {
    "@playwright/test": "^1.49.0",
    "allure-playwright": "^3.0.0",
    "allure-commandline": "^2.27.0",
    "dotenv": "^16.4.7",
    "typescript": "^5.7.0"
  },
  "dependencies": {
    "@axe-core/playwright": "^4.11.2",
    "ws": "^8.20.0"
  }
}
\end{lstlisting}

\section{Configuration JaCoCo (couverture de code Java)}

\begin{lstlisting}[language=XML, caption={pom.xml --- Plugin JaCoCo pour la couverture de code}, label={lst:jacoco}]
<plugin>
  <groupId>org.jacoco</groupId>
  <artifactId>jacoco-maven-plugin</artifactId>
  <version>0.8.12</version>
  <executions>
    <execution>
      <id>prepare-agent</id>
      <goals>
        <goal>prepare-agent</goal>
      </goals>
    </execution>
    <execution>
      <id>report</id>
      <phase>test</phase>
      <goals>
        <goal>report</goal>
      </goals>
    </execution>
  </executions>
</plugin>
\end{lstlisting}

\section{Pipeline GitLab CI}

\begin{lstlisting}[language=bash, caption={.gitlab-ci-tests.yml --- Pipeline de tests (extrait)}, label={lst:gitlab-ci}]
# Integration dans le pipeline principal
include:
  - local: 'tests/ci/.gitlab-ci-tests.yml'

stages:
  - build
  - test          # Nouveau stage pour le robot
  - deploy
  - report        # Nouveau stage pour les rapports

# Tests API (Playwright)
test-robot:api:
  stage: test
  image: mcr.microsoft.com/playwright:v1.49.0-noble
  script:
    - cd tests/playwright
    - npm ci
    - npx playwright test --project=api-tests
  artifacts:
    paths:
      - tests/playwright/allure-results/
    reports:
      junit: tests/playwright/test-results/results.xml

# Tests unitaires backend (JUnit 5 + JaCoCo)
test-robot:backend-unit-tests:
  stage: test
  image: maven:3.9-eclipse-temurin-17
  script:
    - cd backend
    - mvn test --fail-at-end
  artifacts:
    paths:
      - backend/services/*/target/site/jacoco/
    reports:
      junit: backend/services/*/target/surefire-reports/*.xml

# Tests de performance (k6)
test-robot:performance-load:
  stage: test
  image:
    name: grafana/k6:latest
    entrypoint: ['']
  script:
    - k6 run tests/performance/k6-load-test.js
  rules:
    - if: '$CI_COMMIT_BRANCH == "staging"'

# Rapport Allure
test-robot:allure-report:
  stage: report
  script:
    - npm install -g allure-commandline
    - allure generate tests/playwright/allure-results
      --clean -o allure-report
  artifacts:
    paths:
      - allure-report/
    expire_in: 30 days
\end{lstlisting}

\section{Configuration Patrol (tests mobiles Flutter)}

\begin{lstlisting}[language=bash, caption={Configuration Patrol dans pubspec.yaml}, label={lst:patrol}]
# Ajout dans pubspec.yaml de chaque app Flutter
dev_dependencies:
  flutter_test:
    sdk: flutter
  integration_test:
    sdk: flutter
  patrol: ^3.13.0

# Execution des tests
# Standard (sans interactions natives) :
flutter test integration_test/app_test.dart

# Avec Patrol (permissions, dialogs natifs) :
dart pub global activate patrol_cli
patrol test --target integration_test/patrol_test.dart
\end{lstlisting}

\section{Couverture Angular (Istanbul/Karma)}

\begin{lstlisting}[language=Java, caption={angular.json --- Activation de la couverture de code}, label={lst:angular-coverage}]
// Dans angular.json > projects > admin-panel > architect > test
"options": {
  "codeCoverage": true,
  "codeCoverageExclude": [
    "src/test-setup.ts",
    "**/*.spec.ts",
    "**/environments/**"
  ]
}

// Execution : npx ng test --watch=false --browsers=ChromeHeadless
// Resultat : coverage/index.html
\end{lstlisting}
